\documentclass{article}
\usepackage{amsmath}

\usepackage{graphicx}
\usepackage{subcaption}
\usepackage{subcaption}
\usepackage{float}

\title{Mixing Optimization : Optimal Permutation Matrix}


\author{Octave Litrico}
\date{\today}
\begin{document}

\maketitle

\section{Introduction}

I have developed a numerical algorithm to model algal growth in a water basin. 
The model is based on the introduction of a source term into the discretization of the advection--diffusion equation with an upwind discretization of the advection term. 
This additional term modifies the Courant--Friedrichs--Lewy (CFL) condition and, consequently, the value of the time step $\Delta t$. To converge, the numerical scheme must verify : 
\[ \Delta_t \leq \frac{\Delta_x^2}{2\mu + a\Delta_x}\]
Let us first represent the time evolution of this upwind numerical scheme, without any source term, in order to highlight its convergence.

\begin{figure}[H]
    \centering

    \begin{subfigure}{0.45\textwidth}
        \centering
        \includegraphics[width=\textwidth]{C:/Users/octav/OneDrive/modélisation algues/schéma décentré t=0.png}
        \caption{$t = 0\,\mathrm{s}$}
        \label{fig:t0}
    \end{subfigure}
    \hfill
    \begin{subfigure}{0.45\textwidth}
        \centering
        \includegraphics[width=\textwidth]{C:/Users/octav/OneDrive/modélisation algues/schéma décentré t=10.png}
        \caption{$t = 10\,\mathrm{s}$}
        \label{fig:t10}
    \end{subfigure}
\end{figure}
\begin{figure}[H]
    \begin{subfigure}{0.45\textwidth}
        \centering
        \includegraphics[width=\textwidth]{C:/Users/octav/OneDrive/modélisation algues/schéma décentré t=20.png}
        \caption{$t = 20\,\mathrm{s}$}
        \label{fig:t20}
    \end{subfigure}
    \begin{subfigure}{0.45\textwidth}
        \centering
        \includegraphics[width=\textwidth]{C:/Users/octav/OneDrive/modélisation algues/schéma décentré t=30.png}
        \caption{$t = 30\,\mathrm{s}$}
        \label{fig:t30}
    \end{subfigure}

    \caption{Time evolution of the upwind numerical scheme.}
    \label{fig:time_evolution}
\end{figure}
\section{introduction of a biomass source term}

In this report, we present first numerical results illustrating the total biomass in the basin after a growth period of $T = 100\,\mathrm{s}$. 
The algae follow the Han growth model, which is of Haldane if the dynamic of the open state A is at its steady state.
The light intensity follows a Beer-lambert law. The $\epsilon$ light extinction coefficient will be computed once and for all based on an approximation of the affine model : 
\[\epsilon(\rho) = \alpha_0\rho + \alpha_1\]
Where $\rho$ is the biomass concentration, $\alpha_0 > 0 \ [m^2g^{-1}]$ stands for the specific light extinction coefficient of the microalgae species
and $\alpha_1 \ [m^{-1}]$ defines the background turbidity due to all non-microalgae components. We choose the following values : \\
$\alpha_0 = 0.214$ (Monodus subterraneu)
$\alpha_1 = 0.2$ (water extinction coefficient). \\

Therefore, we approximate $\epsilon$ to $\frac{1}{H}\int_{0}^{H}(\alpha_0\rho(z) + \alpha_1)dz$. Let us once again simplify this formula by using a uniform initial concentration :
$\rho(z) = \rho = \frac{B_{tot}}{V} = 1$ with $B_{tot}$ the total initial biomass and $V$ the basin volume. 
We have $\epsilon = 0.207$.
And for the values of the constants involved in the computing of $\mu$ : \\
\[
\begin{aligned}
I_s &= 1500 \ \mathrm{m^{-2}\,\mu mol\,s^{-1}},\\
k_d &= 2.99 \times 10^{-5},\\
k_r &= 4.8 \times 10^{-5} \ \mathrm{s^{-1}},\\
k_h &= 3.64 \times 10^{-5},\\
\tau &= 6.849 \ \mathrm{s},\\
\sigma_H &= 2.9 \times 10^{-4} \ \mathrm{m^{2}\,\mu mol^{-1}}
\end{aligned}
\]

We have 
\[
A = \frac{1}{\frac{k_d}{k_r}\,\tau\,(\sigma_H I)^2 + \tau\,\sigma_H I + 1}
\]

And 

\[
\mu = A\, I\, k_h\, \sigma_H = 1.78\cdot 10^{-4}
\] 
The simulations are performed for a single layer, and therefore without any mixing effects. 
The only varying parameter is the water velocity $u$. For a discretization of 50 points, we impose an initial sinusoidal concentration and an initial total biomass of 42.3 Kg.
After an hour of growth, we get the following total biomass table for different current speed values.

\begin{table}[h!]
\centering
\begin{tabular}{|c|c|}
\hline
\textbf{Current speed (m.s$^{-1}$)} & \textbf{Total biomass (kg)} \\
\hline
0   & 44.5209 \\
0.1 & 44.5208 \\
0.2 & 44.5207 \\
0.3 & 44.5206 \\
0.5 & 44.5205 \\
1   & 44.5204 \\
\hline
\end{tabular}
\caption{Total biomass as a function of current speed}
\label{tab:biomass_speed}
\end{table}

It seems, that the biomass grows less as the current increases. One must take into account that $\Delta_t$, the time step of the numerical scheme depends on the value of the speed $u$. 
This follows from the upwind scheme CFL condition. 


\section{Optimization of the Mixing Matrix}

\subsection{General Principle of the Algorithm}

The objective of the algorithm is to optimize a mixing matrix applied to a discretized basin,
in order to maximize the total biomass produced by the system.

For a given initial configuration, a candidate mixing matrix is repeatedly applied to the system.
The temporal evolution of the basin is simulated with a constant time step, for a given period $T$.
The total biomass is then evaluated.

The algorithm explores different admissible mixing matrices (among every existing permutation matrices) and retains the one
that maximizes the final total biomass.

\subsection{Tested Initial Configurations}

In this study, the basin consists of three vertical layers.
The initial conditions are constructed by placing all the initial biomass (50 Kg) in a single layer,
while the other layers start with zero biomass. 
Three configurations are therefore considered, corresponding to the successive activation of each layer. The idea is to choose the mixing matrix which induces a maximal biomass after an hour of growth.

\subsection{Optimization Results}

For each initial configuration, the optimization algorithm converges to a matrix
that maximizes the final total biomass.

\subsubsection{Configuration 1: biomass initially in layer 1}

The maximal total biomass obtained is:
\[
B_{\text{tot}} = 52.76 \ \mathrm{Kg}
\]

The optimal mixing matrix is:
\[
\mathbf{M}_2 =
\begin{pmatrix}
0 & 0 & 1 \\
0 & 1 & 0 \\
1 & 0 & 0
\end{pmatrix}
\]

This matrix induces a permutation of layers which favorises a complete vertical redistribution of biomass, and in particular, the allocation of layer 1 in position 3.

\subsubsection{Configuration 2: biomass initially in layer 2}

The maximal total biomass obtained is:
\[
B_{\text{tot}} = 52.83 \ \mathrm{Kg}
\]

The optimal mixing matrix is:
\[
\mathbf{M}_3 =
\begin{pmatrix}
1 & 0 & 0 \\
0 & 0 & 1 \\
0 & 1 & 0
\end{pmatrix}
\]

Here, a partial permutation of layers (2 and 3) appears as the optimal strategy.

\subsubsection{Configuration 3: biomass initially in layer 3}

The maximal total biomass obtained is:
\[
B_{\text{tot}} = 52.90 \ \mathrm{Kg}
\]

The optimal mixing matrix is:
\[
\mathbf{M}_1 =
\begin{pmatrix}
1 & 0 & 0 \\
0 & 1 & 0 \\
0 & 0 & 1
\end{pmatrix}
\]

This configuration corresponds to no mixing between layers,
indicating that maintaining the biomass in layer 3 is optimal.


\subsubsection{Comparative Analysis}

These results show that the optimal mixing matrix depends strongly on the initial configuration of the basin.
When biomass is localized in an inherently favorable layer,
no mixing is optimal.
Conversely, in less favorable configurations,
vertical mixing allows biomass to be redistributed to more productive layers. We see that the optimal matrices allocate biomass to the last layer (the deepest one).
Indeed, the $\mu$ values of each layer are the following :

\begin{table}[H]
\centering
\begin{tabular}{|c|c|}
\hline
\textbf{Layer} & \textbf{$\mu$ (s$^{-1}$)} \\
\hline
Layer 1 & $1.4198 \times 10^{-5}$ \\
Layer 2 & $1.4883 \times 10^{-5}$ \\
Layer 3 & $1.5586 \times 10^{-5}$ \\
\hline
\end{tabular}
\caption{Growth rates $\mu$ for each initial layer configuration}
\label{tab:mu_layers}
\end{table}

Therefore the last layer has the biggest growth rate.

\subsection{other initial configurations, 5 layers}

\subsubsection{Configuration 1: biomass initially in layer 1}
After 750 $s$ of growth, the total biomass obtained is:
\[
B_{\text{tot}} = 50.571 \ \mathrm{Kg}
\]

Optimal mixing matrix:
\[
\mathbf{M}_1 =
\begin{pmatrix}
0 & 0 & 0 & 1 & 0 \\
0 & 1 & 0 & 0 & 0 \\
0 & 0 & 1 & 0 & 0 \\
0 & 0 & 0 & 0 & 1 \\
1 & 0 & 0 & 0 & 0
\end{pmatrix}
\]

\subsubsection{Configuration 2: biomass initially in layer 2}

Total biomass obtained:
\[
B_{\text{tot}} = 50.577 \ \mathrm{Kg}
\]

Optimal mixing matrix:
\[
\mathbf{M}_2 =
\begin{pmatrix}
1 & 0 & 0 & 0 & 0 \\
0 & 0 & 0 & 1 & 0 \\
0 & 0 & 1 & 0 & 0 \\
0 & 0 & 0 & 0 & 1 \\
0 & 1 & 0 & 0 & 0
\end{pmatrix}
\]

\subsubsection{Configuration 3: biomass initially in layer 3}

Total biomass obtained:
\[
B_{\text{tot}} = 50.584 \ \mathrm{Kg}
\]

Optimal mixing matrix:
\[
\mathbf{M}_3 =
\begin{pmatrix}
1 & 0 & 0 & 0 & 0 \\
0 & 1 & 0 & 0 & 0 \\
0 & 0 & 0 & 1 & 0 \\
0 & 0 & 0 & 0 & 1 \\
0 & 0 & 1 & 0 & 0
\end{pmatrix}
\]

\subsubsection{Configuration 4: biomass initially in layer 4}

Total biomass obtained:
\[
B_{\text{tot}} = 50.594 \ \mathrm{Kg}
\]

Optimal mixing matrix:
\[
\mathbf{M}_4 =
\begin{pmatrix}
0 & 1 & 0 & 0 & 0 \\
1 & 0 & 0 & 0 & 0 \\
0 & 0 & 1 & 0 & 0 \\
0 & 0 & 0 & 0 & 1 \\
0 & 0 & 0 & 1 & 0
\end{pmatrix}
\]

\subsubsection{Configuration 5: biomass initially in layer 5}

Total biomass obtained:
\[
B_{\text{tot}} = 50.604 \ \mathrm{Kg}
\]

Optimal mixing matrix:
\[
\mathbf{M}_5 =
\begin{pmatrix}
1 & 0 & 0 & 0 & 0 \\
0 & 1 & 0 & 0 & 0 \\
0 & 0 & 1 & 0 & 0 \\
0 & 0 & 0 & 1 & 0 \\
0 & 0 & 0 & 0 & 1
\end{pmatrix}
\]
\end{document}